\section{Cross-Cultural Stress Tests}\label{sec:cross-cultural}

This section records only source-controlled comparisons.  It distinguishes
whole-narrative equivalence from equivalence of a deliberately projected
substructure.  A projected isomorphism is not a claim that two myths have the
same meaning, origin, or historical transmission.

\subsection{Inana: three-night threshold and seven-gate completion}

The Sumerian composition \emph{Inana's Descent to the Nether World} gives a
particularly strong test because both the numerical marker and the state
change are explicit.  Inana is turned into a corpse and hung on a hook.  The
text then states that after three days and three nights have passed,
Nin\v{s}ubur initiates the prescribed rescue sequence.  Enki eventually sends
the life-giving plant and life-giving water; after these are applied the text
states that Inana rises and then begins her ascent from the underworld.

The conservative graph is therefore
\[
\mathrm{life}\to\mathrm{descent}\to\mathrm{death/liminality}
\xrightarrow[\text{rescue activated}]{3\ \mathrm{days}+3\ \mathrm{nights}}
\mathrm{revival}\to\mathrm{ascent}.
\]
The text does \emph{not} state that revival itself occurs exactly at the end of
the third day.  Thus comparison with a ``raised on the third day'' formula is
classified as an order-homomorphism of the projected death--liminality--life
motif, not an exact temporal isomorphism.

The same composition describes seven gates.  At the sixth gate Inana loses
her measuring rod and line; at the seventh she loses the final garment of
ladyship.  Immediately after the completion of this stripping sequence she is
judged and becomes a corpse.  This yields a projected boundary geometry
\[
6\longrightarrow7\longrightarrow\text{regime change},
\]
but with a boundary semantic of death rather than Sabbath cessation.
Consequently it is an order-homomorphism to the Genesis six-to-seven boundary,
not a semantic isomorphism.

\subsection{Gilgamesh XI: a six-to-seven cessation schedule}

The Standard Babylonian flood account states that the storm/flood acts for six
days and seven nights and that, when the seventh day arrives, the gale relents
and the Deluge ends.  After projection to numbered active units and the
cessation boundary, the schedule is
\[
(1,2,3,4,5,6)_{\mathrm{active}}\to7_{\mathrm{cessation}}.
\]
This is isomorphic to the \emph{projected schedule} of six active creation days
followed by Sabbath cessation.  It is explicitly \emph{not} a claim of
whole-narrative identity, common theology, or statistical independence; the
ancient Near Eastern texts belong to historically connected cultural spaces.

\subsection{Zoroastrian three-to-four post-mortem transition}

Zoroastrian sources preserve a three-night liminal interval after death.  A
standard scholarly synthesis of the V\={i}d\={e}vd\={a}d and later tradition
places the soul near the corpse for three days/nights, with movement toward the
\v{C}inwad bridge and judgment at the dawn after the third day.  The shape is
therefore
\[
\mathrm{death}\to[3\ \mathrm{nights}]\to\mathrm{bridge/judgment},
\]
which supports a three-interval threshold motif but not resurrection.

\subsection{Katha Upanishad: three nights in the house of Death}

In \emph{Katha Upanishad} 1.1.7--9, Naciketas remains for three nights fasting
in the house of Yama.  Yama then returns and grants three boons, one for each
night.  This is a clean primary-text instance of a three-night liminal stay
followed by a change of regime:
\[
\mathrm{arrival\ at\ Death}\to[3\ \mathrm{nights}]\to
\mathrm{encounter/boons/teaching}.
\]
It is an analogy to the threshold operator, not a resurrection isomorphism.

\subsection{Kojiki opening: a real \texorpdfstring{$3\mid2$}{3|2} prefix}

The opening of the \emph{Kojiki} groups the first three deities as solitary
deities who appear and hide their bodies.  It then gives two further solitary
deities that likewise appear and hide their bodies.  The source therefore
contains an explicit
\[
3\mid2
\]
prefix.  No primary-text basis has yet been established for continuing this as
$3\mid2\mid1\to0$.  It is retained as a prefix analogy and a target for future
falsification, not as a completed match.

\subsection{Negative controls: Iranian creation counts and Enuma Eli\v{s}}

Middle Persian Zoroastrian cosmogony preserves both six-creation and
seven-stage variants.  In one six-creation account the six g\={a}h\={a}nb\={a}rs
end with man; other accounts enumerate seven creations, for example by adding
fire.  Because the seventh item is not consistently a cessation state, this is
not evidence for the operator $6\to7_{\mathrm{postclosure}}$; the count alone
is a number match.

Likewise, \emph{En\=uma Eli\v{s}} places the creation of humanity in Tablet VI
and devotes Tablet VII predominantly to Marduk's names and praise.  However,
the name-proclamation already begins in the latter part of Tablet VI.  Hence a
clean ``six creation tablets then seventh post-creation tablet'' boundary would
be an overstatement.  This case is retained as a negative control against
retrospective boundary sharpening.

\subsection{Current relation matrix}

\begin{center}
\begin{tabular}{p{0.26\linewidth}p{0.25\linewidth}p{0.28\linewidth}}
\textbf{Comparison} & \textbf{Class} & \textbf{Scope}\\\hline
Genesis / Gilgamesh XI & projected isomorphism & $6\to7$ active-to-cessation schedule\\
Genesis / Inana gates & order-homomorphism & $6\to7$ boundary geometry, different semantics\\
Third-day resurrection / Inana & order-homomorphism & death--liminality--life projection\\
Third-day motif / Zoroastrian soul & analogy & three-night liminality then judgment\\
Third-day motif / Naciketas & analogy & three-night liminality then encounter\\
$3\mid2\mid1\to0$ / Kojiki & analogy & explicit $3\mid2$ prefix only\\
$6\to7$ / Iranian creation variants & number match only & no stable cessation semantics\\
$6\to7$ / En\=uma Eli\v{s} & analogy / negative control & boundary crosses Tablets VI and VII\\
\end{tabular}
\end{center}

No verified cross-cultural full-narrative isomorphism of
$3\mid2\mid1\to0$ is claimed at this stage.
